\paragraph*{04.05.07.01 Projektkosten abschätzen}
\kcilabel{para:04.05.07.01_Projektkosten}{04.05.07.01}
\label{para:04.05.07.01_Projektkosten}

% Task
\par Die Kostenerwartung für das \gls{wsr} und \gls{pb} lag bei ca.~30\,\% der Gesamtprogrammkosten und stellte damit den höchsten Einzelposten dar. Die Kostenschätzung erfolgte auf Basis von Erfahrungswerten aus vergleichbaren Transformationsprojekten wie \gls{vita} und \gls{adam20}. Meine Aufgabe bestand darin, die Kosten für die Entwicklung und Implementierung des \gls{wsr} und \gls{pb} je Welle abzuschätzen, um eine belastbare Grundlage für die Budgetplanung und -kontrolle zu schaffen.

% Action
\par Zur Abschätzung der Projektkosten in den jeweiligen Wellen verwendete ich eine analogiebasierte Top-Down-Betrachtung. Dabei berücksichtigte ich verschiedene Einflussfaktoren, darunter die Komplexität der Entwicklungen, die Anzahl der voraussichtlich beteiligten Ressourcen, die Dauer der Entwicklungsphasen sowie die spezifischen fachlichen und technischen Anforderungen des \gls{wsr} und \gls{pb}. Basale Kosten wie Infrastruktur-, Lizenz- und Betriebskosten wurden bewusst nicht in die Schätzung einbezogen, da diese bereits in anderen Bereichen des Programms budgetiert waren. Darüber hinaus wurde berücksichtigt, dass grundlegende Funktionalitäten wie der Zahl- und Mahnlauf bereits in der ersten Welle implementiert werden mussten, während in den Folgewellen weitere \gls{bk}-spezifische Funktionalitäten ergänzt wurden, was sich auf die Kostenverteilung auswirkte.

% Result
\par Die Kostenschätzungen für die Entwicklung und Implementierung des \gls{wsr} und \gls{pb} wurden termingerecht je Welle erstellt und als Entscheidungsgrundlage für die Budgetfreigabe verwendet. Dadurch konnte die Finanzplanung strukturiert entlang der Umsetzungswellen aufgesetzt und nachvollziehbar fortgeschrieben werden. Mit der Aufnahme von \gls{td} infolge der Rückstellung der \gls{aif}-Implementierung erhöhte sich jedoch die Unsicherheit einzelner Annahmen, sodass die Schätzung als belastbare Planungsbasis mit explizitem Aktualisierungsbedarf zu bewerten war.

% Reflect
\par Die Kostenschätzung für komplexe Transformationsvorhaben wie die Entwicklung und Implementierung des \gls{wsr} und \gls{pb} ist mit erheblichen Unsicherheiten verbunden, insbesondere bei hoher Scope-Dynamik und einer unterkomplexen technischen Voranalyse (vgl. \autoref{fig:R2HErgebnisdokumentationStreamI2C}). Für künftige Wellen wird daher eine frühere Validierung technischer Annahmen, eine transparente Dokumentation der Schätzprämissen sowie eine regelmäßige Nachschätzung bei Scope-Änderungen konsequent eingeplant.